\section{Окружения пакета amsmath}
\par В этом разделе мы рассмотрим окружения пакета \verb"amsmath", предназначенные для вёрстки многострочных формул. Эти окружения очень важны при наборе сложных математических выражений, систем уравнений и цепочек преобразований.

\subsection{Окружение equation и split}
\par Окружение \verb"equation" используется для создания нумерованных уравнений.\index{нумерованное уравнение=equation} Внутри него можно использовать окружение \verb"split", чтобы разбить длинное выражение на несколько строк, сохраняя единый номер.\index{разбиение уравнения на строки=split} Разделителем строк служит \verb"\\", а символ \verb"&" обозначает точку выравнивания:
\begin{equation}\label{eq:one}
    \begin{split}
      a & = b+c-d\\
        & \quad +e-f\\
        & =g+h\\
        & =i
    \end{split}
\end{equation}

\subsection{Окружение multline}
\par Окружение \verb"multline" предназначено для уравнений, не помещающихся в одну строку.\index{многострочное уравнение=multline} Отличие от \verb"split" в том, что выравнивание происходит автоматически: первая строка прижимается влево, последняя --- вправо, а промежуточные строки центрируются:
\begin{multline}
    a+b+c+d+e+f\\
    +i+j+k+l+m+n
\end{multline}

\subsection{Окружение gather}
\par Окружение \verb"gather" центрирует несколько формул по вертикали, не выравнивая их горизонтально друг относительно друга.\index{группировка уравнений=gather} Каждая формула получает свой номер:
\begin{gather}
    a_1=b_1+c_1\\
    a_2=b_2+c_2-d_2+e_2
\end{gather}

\subsection{Окружение align}
\par Окружение \verb"align" позволяет выровнять несколько уравнений по определённой точке (обычно по знаку равенства), обозначаемой символом \verb"&".\index{выравнивание уравнений=align} Каждая строка нумеруется отдельно:
\begin{align}
    a_1& =b_1+c_1\\
    a_2& =b_2+c_2-d_2+e_2
\end{align}
Можно задать и несколько точек выравнивания, чтобы создать несколько столбцов с уравнениями:
\begin{align}
    a_1& =b_1+c_1  & c_1 &= r_1\\
    a_2& =b_2+c_2-d_2+e_2  & c_2 &= r_2
\end{align}

\subsection{Окружение alignat}
\par Окружение \verb"alignat" похоже на \verb"align", но даёт более тонкий контроль над горизонтальными отступами между столбцами.\index{выравнивание уравнений без отступов=alignat} Его обязательный аргумент --- количество <<пар столбцов>>. Сравните результат с \verb"align":
\begin{alignat}{2}
    a_1& =b_1+c_1  &\quad c_1 &= r_1\\
    a_2& =b_2+c_2-d_2+e_2  &\quad c_2 &= r_2
\end{alignat}

\section{Ненумерованные варианты}
\par Все перечисленные выше окружения автоматически нумеруют формулы. Если нумерация не нужна, используйте версии со звёздочкой: \verb"equation*", \verb"multline*", \verb"gather*", \verb"align*".\index{ненумерованные уравнения=equation*=gather*=align*=multline*}
\par Впрочем, заметьте, что \verb"equation*" не даёт никакого преимущества перед обычной записью \verb"\[…\]".
\par Если же нужно убрать номер лишь у одного уравнения внутри нумерованного окружения, к этой строке можно добавить команду \verb"\notag" или \verb"\nonumber":\index{подавление нумерации строки=\notag=\nonumber}
\begin{align}
    x &= y+z \notag\\
    a &= b+c
\end{align}

\section{Системы уравнений}
\par Для набора систем уравнений с фигурной скобкой слева используют окружение \verb"cases":\index{система уравнений=cases}
\[
    |x| = \begin{cases}
        x,  & \text{если } x \geqslant 0,\\
        -x, & \text{если } x < 0.
    \end{cases}
\]
Обратите внимание, что пояснительный текст внутри формулы оборачивается в \verb"\text{}", а символ \verb"&" разделяет формулу и условие.

\section{Упражнение: окружения amsmath}
\begin{staticpart}
Наберите формулы, выбирая подходящие окружения.
\subsection{Первое}
\htmlblockquote{
\begin{equation}
    e=mc^2
\end{equation}
}

\subsection{Второе}
\htmlblockquote{
\begin{gather*}
    d = \sqrt{ (x_2 - x_1)^2 + (y_2 - y_1)^2 } \\
    \cos \alpha = \frac{b^2 + c^2 - a^2}{2bc}
\end{gather*}
}

\subsection{Третье}
\htmlblockquote{
\begin{multline*}
    P(x) = a_0 + a_1 x + a_2 x^2 + a_3 x^3 + \cdots + \\
    + a_n x^n = \sum_{k=0}^n a_k x^k
\end{multline*}
}

\subsection{Четвёртое}
\htmlblockquote{
\begin{align}
    \sin 2x &= 2 \sin x \cos x \\
    1 &= \sin ^2 x + \cos ^2 x
\end{align}
}

\subsection{Пятое}
\htmlblockquote{
\[
    f(x) = \begin{cases}
        \sin x,  & \text{если } x \geqslant 0,\\
        x^2 + 1, & \text{если } x < 0.
    \end{cases}
\]
}

\end{staticpart}

Первое:
\[
    e=mc^2
\]

\par Второе:
\[
    d = \sqrt{ (x_2 - x_1)^2 + (y_2 - y_1)^2 }
    \cos \alpha = \frac{b^2 + c^2 - a^2}{2bc}
\]

\par Третье:
\[
    P(x) = a_0 + a_1 x + a_2 x^2 + a_3 x^3 + \cdots +
    + a_n x^n = \sum_{k=0}^n a_k x^k
\]

\par Четвёртое:
\[
    \sin 2x = 2 \sin x \cos x
    1 = \sin ^2 x + \cos ^2 x
\]

\par Пятое:
\[
    f(x) = ?
\]
